% =====================================================================
%  FICHE 8 — THÈME 3 — NIVEAU MOYEN — ÉNONCÉS
% =====================================================================
\documentclass[11pt,a4paper]{article}
\usepackage[utf8]{inputenc}
\usepackage[T1]{fontenc}
\usepackage{lmodern}
\usepackage[french]{babel}
\usepackage{amsmath,amssymb,mathtools}
\usepackage{icomma}
\usepackage{siunitx}
\sisetup{output-decimal-marker={,},per-mode=power,group-separator={\,},group-minimum-digits=5}
\usepackage[version=4]{mhchem}
\usepackage{xcolor}
\usepackage{tikz}
\usepackage[most]{tcolorbox}
\usepackage{enumitem}
\usepackage{array}
\usepackage{fancyhdr}
\usepackage[hidelinks]{hyperref}
\usetikzlibrary{arrows.meta,positioning,calc}
\usepackage[a4paper,margin=2.1cm,headheight=26pt,headsep=10pt]{geometry}
\usepackage{titlesec}

\definecolor{primary}{RGB}{150,98,162}
\definecolor{primarydark}{RGB}{92,52,110}

\tcbset{boxbase/.style={enhanced,breakable,boxrule=0.9pt,arc=2.5pt,left=3mm,right=3mm,top=2.3mm,bottom=2.3mm,fonttitle=\bfseries,coltitle=white,toptitle=0.6mm,bottomtitle=0.6mm}}
\newtcolorbox[auto counter]{exercice}[1][]{boxbase,colback=primary!4,colframe=primary,title={\textsc{Exercice}~\thetcbcounter\ \ifx\relax#1\relax\relax\else\textmd{--- \itshape #1}\fi}}

\pagestyle{fancy}
\fancyhf{}
\fancyhead[L]{\small\textcolor{primarydark}{\textbf{Fiche 8} — Thème 3 : Calorimétrie}}
\fancyhead[R]{\small\textcolor{primarydark}{\textsc{Moyen}}}
\fancyfoot[C]{\small\thepage}
\renewcommand{\headrulewidth}{0.8pt}
\renewcommand{\headrule}{\hbox to\headwidth{\color{primary}\leaders\hrule height \headrulewidth\hfill}}
\setlength{\parindent}{0pt}\setlength{\parskip}{3pt}

\begin{document}

\begin{center}
\begin{tikzpicture}
\node[rectangle,rounded corners=7pt,fill=primary,text=white,minimum width=16cm,
      minimum height=1.1cm,align=center,inner sep=6pt]
  {{\large\bfseries Thème 3 — Calorimétrie}\\[1pt]{\small Niveau \textsc{Moyen} — énoncés}};
\end{tikzpicture}
\end{center}

\textit{Données pour l'eau :} $c_{\text{glace}}=\SI{2.060}{\kilo\joule\per\kelvin\per\kilogram}$,\;
$c_{\text{eau}}=\SI{4.185}{\kilo\joule\per\kelvin\per\kilogram}$,\;
$L_{\text{fus}}=\SI{334}{\kilo\joule\per\kilogram}$. Masse volumique de l'eau : \SI{1}{\gram\per\milli\litre}.

\begin{exercice}[chauffage en plusieurs étapes]
On veut porter $m=\SI{0.10}{\kilogram}$ de glace de \SI{-20}{\celsius} jusqu'à de l'eau liquide à \SI{50}{\celsius}.
\begin{enumerate}[label=\alph*),leftmargin=2em,topsep=2pt,itemsep=2pt]
  \item Décrire les \textbf{trois} étapes du parcours et la formule associée à chacune.
  \item Calculer la chaleur de chaque étape.
  \item En déduire la chaleur totale $Q$ à fournir.
\end{enumerate}
\end{exercice}

\begin{exercice}[neutralisation dans un calorimètre]
Dans un calorimètre adiabatique en verre (dont on admet qu'il ne s'échauffe pas), on mélange
\SI{50}{\milli\litre} de soude \ce{NaOH} à \SI{2.0}{\mol\per\litre} et \SI{50}{\milli\litre} d'acide
chlorhydrique \ce{HCl} de même concentration, tous deux à \SI{20}{\celsius}. Après réaction, la
température atteint \SI{30}{\celsius}.
\begin{enumerate}[label=\alph*),leftmargin=2em,topsep=2pt,itemsep=2pt]
  \item Quelle est la masse totale de solution (assimilée à de l'eau) et la variation de température ?
  \item Calculer la chaleur $Q$ captée par l'eau (en \si{\kilo\joule}, puis en \si{\joule}).
  \item La réaction $\ce{HCl}+\ce{NaOH}\rightarrow\ce{NaCl}+\ce{H2O}$ est-elle exo ou endothermique ?
\end{enumerate}
\end{exercice}

\begin{exercice}[température finale d'un mélange]
Dans un calorimètre de capacité négligeable, on mélange \SI{200}{\gram} d'eau à \SI{80}{\celsius} et
\SI{300}{\gram} d'eau à \SI{20}{\celsius}.
\begin{enumerate}[label=\alph*),leftmargin=2em,topsep=2pt,itemsep=2pt]
  \item Écrire le bilan $\sum Q_i = 0$ pour les deux masses d'eau (la capacité $c_{\text{eau}}$ se simplifie).
  \item En déduire la température finale $T_{\text{f}}$ du mélange.
\end{enumerate}
\end{exercice}

\begin{exercice}[enthalpie molaire de dissolution]
Dans un calorimètre de capacité négligeable contenant \SI{200}{\gram} d'eau à \SI{20}{\celsius}, on
dissout \SI{0.50}{\mol} d'un sel. La température monte à \SI{24}{\celsius}.
\begin{enumerate}[label=\alph*),leftmargin=2em,topsep=2pt,itemsep=2pt]
  \item Calculer la chaleur $Q$ captée par l'eau.
  \item La dissolution est-elle exo ou endothermique ? Justifier par le sens de variation de la température.
  \item En déduire l'enthalpie \textbf{molaire} de dissolution $\Delta H$ (avec son signe), en \si{\kilo\joule\per\mol}.
\end{enumerate}
\end{exercice}

\end{document}
